\documentclass[10pt,a4paper]{amsart}
\usepackage[margin=19mm]{geometry}
\usepackage{amsmath,amssymb,mathtools}
\usepackage[scr=rsfs,cal=euler]{mathalfa}
\usepackage{xfrac}
\usepackage[unicode,hidelinks,bookmarks=false]{hyperref}
\newcommand{\ie}{i.e.\ }
\newcommand{\cf}{cf.\ }
\newcommand{\resp}{respectively\ }
\newcommand{\Subsection}{\subsection}
\providecommand{\nobreakdash}{\nobreak\hbox{-}}
\setlength{\emergencystretch}{2em}
\multlinegap0pt
\usepackage{amssymb}
\usepackage{xcolor}
\usepackage{algorithm}\usepackage{algpseudocode}
\newtheorem{assumption}{Assumption}
\newtheorem{lemma}{Lemma}
\newtheorem{proposition}{Proposition}
\newtheorem{theorem}{Theorem}
\title{Structural Compatibility and Uniform Stability of Temporally Degenerate Parabolic Systems}
\author{Amadou Ciss\'e}
\date{Source: arXiv:2609.03136v1 (CC BY 4.0)}
\begin{document}
\maketitle

\noindent\emph{Source and licence.} Structural Compatibility and Uniform Stability of Temporally Degenerate Parabolic Systems. Version arXiv:2609.03136v1 (CC BY 4.0), distributed by arXiv under the Creative Commons Attribution 4.0 International licence, http://creativecommons.org/licenses/by/4.0/, as recorded in the arXiv licence metadata for this exact version and shown on the abstract page https://arxiv.org/abs/2609.03136v1. Full licence notice: \texttt{licenses/arxiv-2609.03136-CC-BY-4.0.txt}.\par\smallskip

\noindent\textit{Editorial scope.} Counted unit: the article's theorem thm:critical-residual-certificate (source0.tex lines 2207--2270). The article's complete original proof of it is delivered with it (source0.tex lines 2272--2612, one \texttt{proof} environment of 341 lines). No mathematics is written, repaired or re-derived: the statement and the proof are the article's own lines, in the article's own notation, and no retained line is substituted or re-flowed.  Retained uncounted as the source-local context the proof needs: lines 430--458; lines 462--465; lines 1108--1119; lines 1656--1702; lines 1753--1794; lines 1941--1981; lines 2069--2132.  Nothing was omitted from any retained span, so the package carries no omission record.  No retained line contains a citation, so the wrapper carries no bibliography and no bibliography entry.  No retained line contains a figure environment, an image-inclusion command or drawing code, so no image is delivered and the package declares no asset dependency.  Editorial formatting only, in addition: an AMS \texttt{amsart} preamble that restates the article's own declarations and macros used by the retained lines, namely the environments \texttt{assumption}, \texttt{lemma}, \texttt{proposition}, \texttt{theorem} on separate counters, together with the standard packages those lines need.  The retained environments print in this document as Assumption 1, Lemma 1, Proposition 1, Theorem 1 in order of appearance, whereas the article prints the same items with the numbering of its own full document.  The complete native source of the article as distributed in the arXiv e-print for this exact version is bundled unchanged as \texttt{sources/209269/source0.tex}.
\begin{assumption}[Temporal coefficients]
\label{ass:time}

The functions \(\alpha,\beta\in C^{1}((0,T])\)
satisfy
\[
\alpha(t)>0,
\qquad
\beta(t)>0,
\qquad
t\in(0,T].
\]

Moreover,
\(\alpha(0)=0\),
and the ratio
\[
d(t):=\frac{\beta(t)}{\alpha(t)}
\]
satisfies
\begin{equation}
0<d_-\le d(t)\le d_+,
\qquad
t\in(0,T],
\label{eq:d-bounds}
\end{equation}
for some constants \(d_->0\) and \(d_+<\infty\).

\end{assumption}

\begin{assumption}[Principal operator]
\label{ass:A0}
The operator \(A_0:D(A_0)\subset X\rightarrow X\) is densely defined, self-adjoint, strictly positive, and possesses a compact resolvent.
\end{assumption}

\begin{equation}
\dot z
+
d(t)A_{0}z
+
\left(
A_{r}-BK_{1}C
\right)z
=
0,
\label{eq:regular-evolution}
\end{equation}

Projecting the evolution equation
\eqref{eq:regular-evolution}
onto \(X_c\) and \(X_s\) yields
\begin{equation}
\begin{aligned}
\dot z_c
&=
G_{cc}(d(t))z_c
+
G_{cs}z_s,
\\
\dot z_s
&=
G_{sc}z_c
+
G_{ss}(d(t))z_s,
\end{aligned}
\label{eq:critical-residual-block-system}
\end{equation}
where
\begin{equation}
G_{cc}(d)
:=
-dA_{0,c}
-
P_NL_KP_N,
\label{eq:Gcc-definition}
\end{equation}
\begin{equation}
G_{ss}(d)
:=
-dA_{0,s}
-
Q_NL_KQ_N,
\label{eq:Gss-definition}
\end{equation}
and
\begin{equation}
G_{cs}
:=
-P_NL_KQ_N,
\qquad
G_{sc}
:=
-Q_NL_KP_N.
\label{eq:coupling-blocks}
\end{equation}

\begin{lemma}[Residual Dissipation Margin]
\label{lem:uniform-high-frequency-dissipation}

Assume that Assumptions~\ref{ass:time} and~\ref{ass:A0} hold. For a fixed integer \(N\ge1\), define
\begin{equation}
\mu_s
:=
d_-\lambda_{N+1}
-
\ell_s(K_1).
\label{eq:residual-margin}
\end{equation}

Then, for every
\[
d\in[d_-,d_+]
\]
and every
\[
x_s\in D(A_0)\cap X_s,
\]
one has
\begin{equation}
\operatorname{Re}
\left\langle
G_{ss}(d)x_s,x_s
\right\rangle_X
\le
-\mu_s\|x_s\|_X^2.
\label{eq:uniform-residual-dissipation}
\end{equation}

In particular, if
\begin{equation}
d_-\lambda_{N+1}
>
\ell_s(K_1),
\label{eq:positive-residual-margin}
\end{equation}
then the residual dynamics are uniformly dissipative over \(d\in[d_-,d_+]\).

\end{lemma}

\begin{lemma}[Weighted Spillover Estimates]
\label{lem:weighted-spillover-estimates}

For every
\[
x_c\in X_c,
\qquad
x_s\in X_s,
\]
one has
\begin{equation}
\left|
\left\langle
\Pi G_{cs}x_s,x_c
\right\rangle_X
\right|
\le
\gamma_{cs}
\,
\|\Pi^{1/2}x_c\|_X
\,
\|x_s\|_X,
\label{eq:weighted-critical-to-residual-estimate}
\end{equation}
and
\begin{equation}
\left|
\left\langle
G_{sc}x_c,x_s
\right\rangle_X
\right|
\le
\gamma_{sc}
\,
\|\Pi^{1/2}x_c\|_X
\,
\|x_s\|_X.
\label{eq:weighted-residual-to-critical-estimate}
\end{equation}

\end{lemma}

\begin{proposition}[Common Quadratic Certificate]
\label{prop:critical-quadratic-certificate}

Consider the critical subsystem
\begin{equation}
\dot z_c(t)
=
G_{cc}(d(t))z_c(t),
\label{eq:critical-subsystem}
\end{equation}
where
\[
G_{cc}(d)
=
-dA_{0,c}
-
P_NL_KP_N,
\qquad
d\in[d_-,d_+].
\]

Assume that there exist a self-adjoint positive-definite operator
\[
\Pi:X_c\rightarrow X_c
\]
and a constant
\[
\mu_c>0
\]
such that
\begin{equation}
\Pi G_{cc}(d)
+
G_{cc}(d)^*\Pi
\le
-2\mu_c\Pi,
\qquad
\forall d\in[d_-,d_+].
\label{eq:critical-LMI}
\end{equation}

Then the critical evolution family is uniformly exponentially stable.
More precisely, there exists
\[
M_c
=
\sqrt{
\frac{\lambda_{\max}(\Pi)}
{\lambda_{\min}(\Pi)}
}
\]
such that
\begin{equation}
\|z_c(t)\|_X
\le
M_c
e^{-\mu_c(t-s)}
\|z_c(s)\|_X,
\qquad
0\le s\le t.
\label{eq:critical-exp-estimate}
\end{equation}

\end{proposition}

\begin{theorem}[Uniform Stability Certificate]
\label{thm:critical-residual-certificate}

Assume Assumptions~\ref{ass:time} and~\ref{ass:A0}.

Suppose that

\begin{enumerate}

\item
the residual dissipation margin satisfies
\(
\mu_s>0;
\)

\item
the critical subsystem satisfies the quadratic certificate of
Proposition~\ref{prop:critical-quadratic-certificate};

\item
there exists
\(
\eta>0
\)
such that
\begin{equation}
4\eta\mu_c\mu_s
>
\left(
\gamma_{cs}
+
\eta\gamma_{sc}
\right)^2.
\label{eq:small-gain-condition}
\end{equation}

\end{enumerate}

Then there exist constants
\[
M\ge1,
\qquad
\omega>0,
\]
independent of the admissible trajectory
\[
d(\cdot):[0,\infty)\rightarrow[d_-,d_+],
\]
such that every solution of
\eqref{eq:regular-evolution}
satisfies
\begin{equation}
\|z(t)\|_X
\le
Me^{-\omega(t-s)}
\|z(s)\|_X,
\qquad
0\le s\le t.
\label{eq:uniform-global-decay}
\end{equation}

Hence, the regular evolution family is uniformly exponentially stable.

\end{theorem}

\begin{proof}

Consider the block-diagonal Lyapunov functional
\begin{equation}
V(z)
=
\left\langle
\Pi z_c,
z_c
\right\rangle_X
+
\eta\|z_s\|_X^2,
\label{eq:coupled-Lyapunov}
\end{equation}
where \(\eta>0\) satisfies
\eqref{eq:small-gain-condition}.

Since
\[
\Pi=\Pi^*>0
\]
on the finite-dimensional space \(X_c\), there exist constants
\[
m_\Pi
=
\lambda_{\min}(\Pi)>0,
\qquad
M_\Pi
=
\lambda_{\max}(\Pi)>0,
\]
such that
\[
m_\Pi\|z_c\|_X^2
\le
\left\langle
\Pi z_c,
z_c
\right\rangle_X
\le
M_\Pi\|z_c\|_X^2.
\]
Consequently,
\begin{equation}
m_V\|z\|_X^2
\le
V(z)
\le
M_V\|z\|_X^2,
\label{eq:Lyapunov-equivalence}
\end{equation}
where
\[
m_V:=\min\{m_\Pi,\eta\},
\qquad
M_V:=\max\{M_\Pi,\eta\}.
\]

Along a strong solution of
\eqref{eq:critical-residual-block-system}, one has
\[
\begin{aligned}
\dot V
={}&
\left\langle
\left(
\Pi G_{cc}(d(t))
+
G_{cc}(d(t))^*\Pi
\right)z_c,
z_c
\right\rangle_X
%\\
%&
+
2\operatorname{Re}
\left\langle
\Pi G_{cs}z_s,
z_c
\right\rangle_X
%\\
%&
+
2\eta\operatorname{Re}
\left\langle
G_{sc}z_c,
z_s
\right\rangle_X
\\
&+
2\eta\operatorname{Re}
\left\langle
G_{ss}(d(t))z_s,
z_s
\right\rangle_X .
\end{aligned}
\]

By the common quadratic certificate of Proposition~\ref{prop:critical-quadratic-certificate},
\[
\left\langle
\left(
\Pi G_{cc}(d(t))
+
G_{cc}(d(t))^*\Pi
\right)z_c,
z_c
\right\rangle_X
\le
-2\mu_c
\|\Pi^{1/2}z_c\|_X^2.
\]

By Lemma~\ref{lem:uniform-high-frequency-dissipation},
\[
\operatorname{Re}
\left\langle
G_{ss}(d(t))z_s,
z_s
\right\rangle_X
\le
-\mu_s\|z_s\|_X^2.
\]

Finally, the weighted spillover estimates of Lemma~\ref{lem:weighted-spillover-estimates}
give
\[
\left|
\left\langle
\Pi G_{cs}z_s,
z_c
\right\rangle_X
\right|
\le
\gamma_{cs}
\|\Pi^{1/2}z_c\|_X
\|z_s\|_X
\]
and
\[
\left|
\left\langle
G_{sc}z_c,
z_s
\right\rangle_X
\right|
\le
\gamma_{sc}
\|\Pi^{1/2}z_c\|_X
\|z_s\|_X.
\]

Combining these estimates yields
\begin{equation}
\begin{aligned}
\dot V
\le{}&
-2\mu_c
\|\Pi^{1/2}z_c\|_X^2
-
2\eta\mu_s
\|z_s\|_X^2
%\\
%&
+
2\left(
\gamma_{cs}
+
\eta\gamma_{sc}
\right)
\|\Pi^{1/2}z_c\|_X
\|z_s\|_X.
\end{aligned}
\label{eq:coupled-Lyapunov-derivative}
\end{equation}

Introduce
\[
u:=\|\Pi^{1/2}z_c\|_X,
\qquad
v:=\|z_s\|_X,
\]
and define
\[
H_\eta
:=
\begin{pmatrix}
2\mu_c
&
-\left(
\gamma_{cs}+\eta\gamma_{sc}
\right)
\\[1mm]
-\left(
\gamma_{cs}+\eta\gamma_{sc}
\right)
&
2\eta\mu_s
\end{pmatrix}.
\]
Then
\[
\dot V
\le
-
\begin{bmatrix}
u\\
v
\end{bmatrix}^{\!T}
H_\eta
\begin{bmatrix}
u\\
v
\end{bmatrix}.
\]

The first leading principal minor of \(H_\eta\) is positive because
\[
2\mu_c>0,
\]
and its determinant satisfies
\[
\det H_\eta
=
4\eta\mu_c\mu_s
-
\left(
\gamma_{cs}
+
\eta\gamma_{sc}
\right)^2
>
0
\]
by \eqref{eq:small-gain-condition}.
Hence,
\[
H_\eta>0.
\]

Let
\[
D_\eta
:=
\begin{pmatrix}
1&0\\
0&\eta
\end{pmatrix}.
\]
Since both \(H_\eta\) and \(D_\eta\) are positive definite, the constant
\[
\nu
:=
\lambda_{\min}
\left(
D_\eta^{-1/2}
H_\eta
D_\eta^{-1/2}
\right)
\]
is strictly positive. Therefore,
\[
\begin{bmatrix}
u\\
v
\end{bmatrix}^{\!T}
H_\eta
\begin{bmatrix}
u\\
v
\end{bmatrix}
\ge
\nu
\left(
u^2+\eta v^2
\right).
\]
Since
\[
u^2+\eta v^2
=
\left\langle
\Pi z_c,
z_c
\right\rangle_X
+
\eta\|z_s\|_X^2
=
V(z),
\]
it follows that
\begin{equation}
\dot V
\le
-\nu V.
\label{eq:global-Lyapunov-decay}
\end{equation}

Gronwall's inequality gives
\[
V(z(t))
\le
e^{-\nu(t-s)}
V(z(s)),
\qquad
0\le s\le t.
\]
Using
\eqref{eq:Lyapunov-equivalence}, we obtain
\[
\|z(t)\|_X^2
\le
\frac{M_V}{m_V}
e^{-\nu(t-s)}
\|z(s)\|_X^2.
\]
Consequently,
\[
\|z(t)\|_X
\le
\sqrt{\frac{M_V}{m_V}}
e^{-\frac{\nu}{2}(t-s)}
\|z(s)\|_X.
\]

Thus, \eqref{eq:uniform-global-decay}
holds with
\[
M
=
\sqrt{\frac{M_V}{m_V}},
\qquad
\omega
=
\frac{\nu}{2},
\]
and these constants are independent of the admissible trajectory \(d(\cdot)\).

The estimate is first obtained for strong solutions and extends to mild solutions by the standard density argument associated with the well-posed evolution family.

\end{proof}

\end{document}
